\par\noindent\textbf{(1)} 平均场每自旋自由能 f(m)=Jz m²/2−k\_BT ln[2cosh((Jz m+h)/(k\_BT))]，第一项修正平均场重复计数。极值f'=Jz[m−tanh((Jz m+h)/(k\_BT))]=0，故m=tanh[(Jzm+h)/(k\_BT)]，它是自洽条件而不是直接代入外场的单自旋结果。
\par\noindent\textbf{(2)} 零场在m≈0展开tanh x=x−x³/3+…，令Tc=Jz/k\_B，得到0≈(Tc/T−1)m−(Tc/T)\textasciicircum{}3m³/3。T略低Tc时自发磁化m²≈3(Tc−T)/Tc，指数β=1/2；高温m=0稳定，低温出现两个相反符号的对称极值。
\par\noindent\textbf{(3)} 对自洽式求h导，在m=0且T>Tc处χ=∂m/∂h=[1/(k\_BT)]/[1−Jz/(k\_BT)]=1/[k\_B(T−Tc)]，表现平均场临界发散。若观测总磁矩，需再乘每自旋磁矩量纲，不能混同无量纲m的χ与SI磁化率。
\par\noindent\textbf{(4)} 在T=Tc将自洽展开为m≈m+h/(k\_BTc)−(m+h/(k\_BTc))³/3，主阶h/(k\_BTc)≈m³/3，因此m∝h\textasciicircum{}\{1/3\}，对应δ=3。有限N会使真实配分函数解析、严格奇点消失；低维涨落可能使平均场临界指数失效。
